\documentclass{beamer}

\usepackage[utf8]{inputenc}
\usepackage[english]{babel}
\usepackage{textpos}
\usepackage{graphicx}
\usepackage{textcomp}
\usepackage{animate}
\usepackage{comment}
%\usepackage{multicol}
%\setlength{\columnsep}{5cm}
\usepackage{amsmath}
\usepackage{amsfonts}
\usepackage{verbatim}
\usepackage{varwidth}
\usepackage{caption}
\usepackage{dcolumn}
\usepackage{amssymb}
\usepackage{algorithm}
\usepackage{algorithmic}
%\usepackage{physics}
\usepackage{hyperref}
\usepackage{scalerel,stackengine}
%\usepackage{fancyvrb}


\newcommand{\code}[1]{\texttt{#1}}
\newcommand{\shp}{$>>>$ }
\newcommand{\tab}{\hspace*{22pt}}

\usetheme{Madrid}

\title{NumPy and Matplotlib}
\author{EEC 3150}
\date{}

\institute{Trevecca Nazarene University}
 


%%%%%%%%% BEGIN DOCUMENT %%%%%%%%%
%%%%%%%%% BEGIN DOCUMENT %%%%%%%%%
%%%%%%%%% BEGIN DOCUMENT %%%%%%%%%

\begin{document}

% get title page 
\frame{\titlepage}
 
%%%%%%%%%%
%%%%%%%%%%
\begin{frame}	
	\frametitle{Outline}
	\tableofcontents	
\end{frame}


%%%%%%%%%%
%%%%%%%%%%
\section{NumPy}

\begin{frame}	
	\center{\Huge{NumPy: \\ Numerical Python}}
\end{frame}

\subsection{Intro to NumPy}

\begin{frame}	
	\frametitle{Intro to NumPy}
	
	\begin{block}{NumPy:}
		NumPy is one of the most commonly used third-party libraries in Python \\
		\href{https://numpy.org/doc/stable/}{https://numpy.org/doc/stable/}
		
		\vspace{10pt}
		Key features:
		\begin{itemize}
			\item Extensive library of tools for mathematical computations
			\begin{itemize}
				\item Trigonometric, exponential, and special functions
				\item Statistics and random number generation
				\item Linear algebra and Fourier analysis (and more)
			\end{itemize}
			\item Bread and butter: \code{ndarray} objects
			\begin{itemize}
				\item $n$-dimensional arrays
				\item Lists but infinitely better (how you wish lists worked)
			\end{itemize}
		\end{itemize}
		
		\vspace{10pt}
		Importing numpy: \\
		\code{import numpy as np}
	\end{block}
	
\end{frame}


%%%%%%%%%%
%%%%%%%%%%
\subsection{Creating \code{ndarray}'s}

\begin{frame}	
	\frametitle{Creating \code{ndarray}'s}
	
	\begin{block}{Functions for creating arrays:}
		\begin{scriptsize}
		
%		NumPy provides many functions for creating arrays of different kinds
		\begin{itemize}
			\item \code{np.array(x)} will convert the list \code{x} into an array \\
					The function accepts an additional keyword argument \code{dtype} that allows for the specification of data type the array should have (ex: np.int32, np.float64)
			\item \code{np.zeros(n)} will create an array of length $n$ (int) of all zeros \\
					You can also pass a tuple instead of an integer to create a multi-dimensional array (ex: \code{np.zeros((m,n))} creates and $m\times n$ array)
			\item \code{np.ones(n)} (like above, but all ones)
			\item \code{np.full(n,value)} (like above, but fills array with \code{value})
%			\item \code{np.eye(n)} creates an $n\times n$ identity matrix
			\item \code{np.arange(a,b,step)} same as Python \code{range} but allows floats
			\item \code{np.linspace(a,b,n)} creates an array of \code{n} evenly-spaced points from \code{a} to \code{b}
			\item \code{np.random.rand(n)} creates a length \code{n} array of random numbers b/t 0 and 1 \\
					This uses the \code{np.random} library (more on that later)
		\end{itemize}
		Some additional points:
		\begin{itemize}
			\item The different dimensions of a numpy array are called \emph{axes} (\code{axis} is often a keyword argument to many functions)
			\item Use the \code{shape} attribute to get the length of an array along each axis (ex: \code{x.shape})
		\end{itemize}
		
%		Equivalent: \code{2*np.ones((3,3))} and \code{np.full((3,3),2)}
		\end{scriptsize}
	\end{block}
	
\end{frame}

\begin{frame}	
	\frametitle{Creating \code{ndarray}'s}
	
	\begin{exampleblock}{2 ways of creating a $3\times3$ array of all 5's:}
		\begin{itemize}
			\item np.full((3,3),5)
			\item 5*np.ones((3,3))
		\end{itemize}
	\end{exampleblock}
	
	\begin{exampleblock}{Row vs. column vs. 1-D:}
		These are all different objects:
		\begin{itemize}
			\item np.ones(5) \\
					This is a 1-D array
			\item np.ones((5,1)) \\
					This is a 2-D array (column)
			\item np.ones((1,5)) \\
					This is a 2-D array (row)
		\end{itemize}
	\end{exampleblock}
	
\end{frame}

%%%%%%%%%%
%%%%%%%%%%
\subsection{Indexing \code{ndarray}'s}

\begin{frame}	
	\frametitle{Indexing \code{ndarray}'s}
	
	\begin{block}{Indexing \code{ndarray}'s}
		\begin{scriptsize}
		Standard list indexing:
		\begin{itemize}
			\item Index an element: \code{x[1]}
			\item Get a slice: \code{x[4:12]}
			\item Index from the end: \code{x[-3]}
		\end{itemize}
		
		Numpy also supports \emph{fancy indexing}:
		\begin{itemize}
			\item Index a 2-D array: \code{x[2,5]}
			\item Index an entire row, column respectively: \code{x[2,:]}, \code{x[:,5]}
			\item Index a list of specific elements: \code{x[[1,4,8]]}
			\item Index with a boolean mask: \code{x[[True,False,False,True]]}
			\item Index all elements with values greater than 0: \code{x[x>0]} \\
			\begin{itemize}
				\scriptsize
				\item How it works: \code{x>0} creates a boolean mask used to index the \code{x} array
				\item This works for multi-dimensional arrays too
				\item Limitation: you can't do multiple conditionals in the same fancy index
				% (but you can still accomplish the same functionality using Numpy's functions for logical operations)
			\end{itemize}
			\item Get the indices (not values) where an array is greater than 0: \code{np.where(x>0)}
			\item Indexing with \code{np.newaxis} will add another axis to the array
		\end{itemize}
		In addition to accessing array elements with the above techniques, you can also SET elements with the above techniques
		\end{scriptsize}
	\end{block}
	
\end{frame}

\begin{frame}	
	\frametitle{Indexing \code{ndarray}'s}
	
	\begin{block}{Conditional masking vs. np.where()}
%		Let \code{x} be a numpy array:
		\begin{itemize}
			\item \code{x[x>0]} returns the VALUES where \code{x} is greater than 0
			\item \code{np.where(x>0)} returns the INDICES where \code{x} is greater than 0
		\end{itemize}
	\end{block}
	
	\begin{block}{np.where() on multi-dimensional arrays}
		\begin{itemize}
			\item Running a \code{np.where()} on a 2D (or higher) array will produce a tuple of index arrays
			\item Each tuple corresponds to the indices of the respective axis
			\item Indexing the original array with this tuple of index arrays will produce a 1D array
		\end{itemize}
	\end{block}
	
\end{frame}

\begin{frame}	
	\frametitle{Indexing \code{ndarray}'s}
	
	\begin{exampleblock}{Set negative elements to 0:}
		\begin{itemize}
			\footnotesize
			\item Generate some random data
			\item Set all negative values to 0
		\end{itemize}
	\end{exampleblock}
	
	\begin{exampleblock}{Index with multiple conditions:}
		\begin{itemize}
			\footnotesize
			\item Generate some random data
			\item Using several boolean masks, index the values such that $a<x<b$
			\item There are several ways to do this
			\begin{itemize}
				\footnotesize
				\item Use numpy's \code{np.logical\_and(x,y)} function
				\item Repeatedly index and overwrite the same array with different boolean masks
			\end{itemize}
		\end{itemize}
	\end{exampleblock}
	
\end{frame}

%%%%%%%%%%
%%%%%%%%%%
\subsection{Math operations on \code{ndarray}'s}

\begin{frame}	
	\frametitle{Math operations \code{ndarray}'s}
	
	\begin{block}{Math operations \code{ndarray}'s:}
%		\begin{footnotesize}
		
		Numpy arrays can handle mathematical operations in an intuitive, element-wise way:
		\begin{itemize}
			\item We can perform virtually any mathematical operation on any combination of arrays and scalars (as long as their lengths in each dimension match)
			\item The exception to this is Boolean operations, but Numpy has functions for these anyway
			\item We can also apply virtually any of Numpy's mathematical functions on arrays in an element-wise fashion
			\item Numpy can even add arrays of different sizes in certain contexts
		\end{itemize}
%		\end{footnotesize}
	\end{block}
	
\end{frame}

\begin{frame}	
	\frametitle{Math operations \code{ndarray}'s}
	
	\begin{exampleblock}{Plot some random numbers:}
		\begin{itemize}
			\item Create an array of random numbers b/t $-a$ and $a$
			\item Create and plot a linearly-spaced array of numbers with the above random numbers added as noise
		\end{itemize}
	\end{exampleblock}
	
	\begin{exampleblock}{Plot sin and cos:}
		Plot sin$(x)$ and cos$(x)$ from 0 to $2\pi$
	\end{exampleblock}
	
	\begin{exampleblock}{Plot positive sin$(x)$:}
		\begin{itemize}
			\footnotesize
			\item Create \code{x,y} arrays for plotting sin$(x)$ from 0,$4\pi$
			\item Only plot values where sin$(x)\ge0$
			\item Plot with dots
		\end{itemize}
	\end{exampleblock}
	
\end{frame}


%%%%%%%%%%
%%%%%%%%%%
\subsection{Misc. Array Tools}

\begin{frame}	
	\frametitle{Misc. Array Tools}
	
	\begin{block}{Misc. Array Tools}
		\begin{scriptsize}
		\begin{itemize}
			\item \code{np.reshape(x,(m,n))} reshapes \code{x} into a new shape; also, x.reshape((m,n))
			\item \code{x.flatten()} returns a 1-D reshaping of \code{x} (doesn't overwrite \code{x})
			\item \code{np.flipud(x)} (ud: up-down) flips a multi-dim array vertically
			\item \code{np.fliplr(x)} (lr: left-right) flips a multi-dim array horizontally
			\item \code{np.rot90(x,n)} applies an $n*90$ degree rotation to the array
			\item \code{np.roll(x,n)} applies a cyclic shift by $n$ to the right
			\item \code{np.sort(x,axis=-1)} returns a sorted version of \code{x} (doesn't overwrite \code{x})
			\item \code{np.argsort(x,axis=-1)} returns the indices of \code{x} which sort it
			\item \code{x.T} transposes a 2D array
			\item \code{np.unique(x)} returns a sorted array of all the unique elements of \code{x}
			\item \code{np.append(x,value)} returns an array with \code{value} appended to \code{x} (doesn't overwrite \code{x})
			\item \code{np.insert(x,idx,value)} returns an array with \code{value} inserted to \code{x} at index \code{idx} (doesn't overwrite \code{x})
			\item \code{np.delete(x,idx)} returns an array with \code{value} deleted from \code{x} (doesn't overwrite \code{x})
%			\begin{itemize}
%				\scriptsize
%				\item 
%			\end{itemize}
		\end{itemize}
		\end{scriptsize}
	\end{block}
	
\end{frame}

\begin{frame}	
	\frametitle{Misc. Array Tools}
	
	\begin{exampleblock}{Indices of all unique values:}
		\footnotesize
		\begin{itemize}
			\footnotesize
			\item Create a list of random integers
			\item Get unique values
			\item For each unique value, create of list of indices where the array has that value
			\item Remove all zeros from the array
		\end{itemize}
	\end{exampleblock}
	
	\begin{exampleblock}{Sort data before plot:}
		\footnotesize
		\begin{itemize}
			\footnotesize
			\item Create random $x$ data from 0 to $2\pi$
			\item $y = sin(x)$
			\item Plot and see that the data is out of order
			\item Sort the data so that the plot looks correct (argsort on $x$ and apply to $x$ and $y$)
			\item Plot again (comment out old plot)
		\end{itemize}
	\end{exampleblock}
	
\end{frame}

%%%%%%%%%%
%%%%%%%%%%
\subsection{Math Functions}

\begin{frame}	
	\frametitle{Math Functions}
	
	\begin{block}{Math Functions}
		\begin{scriptsize}
		\begin{itemize}
			\item sum, prod, cumsum, cumprod
			\item \code{np.amin(x,axis=None)} return the minimum of an array or minimum along an axis \\
					Also: amax, argmin, and argmax
			\item degrees, radians (conversion)
			\item Trig: sin, ... , arcsin... , sinh, ... , arcsinh, ...
			\item Exp: exp, log, log2, log10
			\item floor, ceil, lcm, gcd
			\item Complex numbers: \code{x = 1+1j}
			\begin{itemize}
				\scriptsize
				\item \code{x.real}, \code{x.imag}, \code{np.conj(x)}, \code{np.angle(x)}
			\end{itemize}
		\end{itemize}
		
		There are also:
		\begin{itemize}
			\item Set functions: union, intersection, ...
			\item Boolean functions: logical\_and, ... , all, isinf
		\end{itemize}
		\end{scriptsize}
	\end{block}
	
\end{frame}


%%%%%%%%%%
%%%%%%%%%%
\subsection{Random Functions}

\begin{frame}	
	\frametitle{Random Functions}
	
	\begin{block}{Random Functions (\code{np.random}):}
		\begin{scriptsize}
		\href{https://numpy.org/doc/stable/reference/random/legacy.html}{https://numpy.org/doc/stable/reference/random/legacy.html}
		
		\vspace{10PT}
		All functions begin with \code{np.random.}
		\begin{itemize}
			\item \code{seed(n)} set the seed for RNG (useful for reproducible random sequences)
			\item \code{rand(n)} create an array of uniform random numbers b/t 0 and 1 (pass multiple args for multi-dim)
			\item \code{randn(n)} create an array of normally-distributed random numbers (pass multiple args for multi-dim)
			\item \code{randint(low,high,size)} create an array of uniform random integers numbers b/t low (inc) and high (excl)
			\item \code{uniform(low,high,size)} create an array of uniform random numbers b/t low and high
			\item \code{normal(loc,scale,size)} create an array of normally-distributed random numbers with mean \code{loc} and std \code{scale}
			\item \code{multivariate\_normal(mean,cov,size)} create an array of normally-distributed random numbers with mean and cov
			\item Many more distributions...
		\end{itemize}
		\end{scriptsize}
	\end{block}
	
\end{frame}

\begin{frame}	
	\frametitle{Random Functions}
	
	\begin{block}{Random Functions (\code{np.random}):}
		\begin{scriptsize}
		\href{https://numpy.org/doc/stable/reference/random/legacy.html}{https://numpy.org/doc/stable/reference/random/legacy.html}
		
		\vspace{10pt}
		All functions begin with \code{np.random.}
		\begin{itemize}
			\item \code{shuffle(x)} randomly shuffle \code{x} in-place (DOES CHANGE \code{x}!!!)
			\item \code{permutation(x)} randomly shuffle \code{x} (doesn't change \code{x})
			\item \code{choice(arr,size,replace=True,p=None)} return a random sample (w/ or w/out replacement) or \code{arr} of \code{size} with optional weights \code{p}
		\end{itemize}
		\end{scriptsize}
	\end{block}
	
\end{frame}

\begin{frame}	
	\frametitle{Math Functions}
	
	\begin{exampleblock}{Plot a multi-normal distribution:}
		\footnotesize
		\begin{itemize}
			\footnotesize
			\item Generate (n,2) array of multi-normal data with some mean and covariance
			\item Plot the data
		\end{itemize}
	\end{exampleblock}
	
	\begin{exampleblock}{Sample from a probability distribution:}
		\footnotesize
		Use np.cumsum and np.random.rand to sample from a PDF of your choice (the np.random library has a function for this, but doing it manually is a good exercise)
		\begin{itemize}
			\footnotesize
			\item Create function $f$ for probabilities
			\item Create grid of $x$ values (limits and number of points)
			\item Sample $f$ on grid (array of $f$ values)
			\item Normalize and cumsum the array
			\item Generate random numbers and determine which sample they fall into
		\end{itemize}
	\end{exampleblock}
	
	\begin{exampleblock}{Sample from a probability distribution:}
		\footnotesize
		\begin{itemize}
			\footnotesize
			\item Redo previous exercise with \code{np.random.choice}
		\end{itemize}
	\end{exampleblock}
	
\end{frame}


%%%%%%%%%%
%%%%%%%%%%
\subsection{Statistics Functions}

\begin{frame}	
	\frametitle{Statistics Functions}
	
	\begin{block}{Statistics Functions:}
		\begin{scriptsize}
		\href{https://numpy.org/doc/stable/reference/routines.statistics.html}{https://numpy.org/doc/stable/reference/routines.statistics.html}
		\begin{itemize}
			\item average, mean, median, std, var
			\item cov, corrcoef
			\item quantile, percentile (functionally equivalent, different inputs)
			\item histogram, histogram2d, histogramdd
			\item bincount (count the number of occurences of unique integers)
		\end{itemize}
		\end{scriptsize}
	\end{block}
	
\end{frame}

\begin{frame}	
	\frametitle{Statistics Functions}
	
	\begin{exampleblock}{Correlate some variables:}
		\footnotesize
		\begin{itemize}
			\footnotesize
			\item Create some linspace-d x data
			\item Create $y = mx+b$ plus normal noise
			\item Plot x, y
			\item Get corrcoef
			\item Vary strength of noise
		\end{itemize}
	\end{exampleblock}
	
\end{frame}


%%%%%%%%%%
%%%%%%%%%%
\subsection{Linear Algebra Functions}

\begin{frame}	
	\frametitle{Linear Algebra Functions}
	
	\begin{block}{Linear Algebra Functions (\code{np.linalg}):}
		\scriptsize
		
		\href{https://numpy.org/doc/stable/reference/routines.linalg.html}{https://numpy.org/doc/stable/reference/routines.linalg.html}
		
		\vspace{10pt}
		NOTE: some of these use \code{np} and some use \code{np.linalg}
		\begin{itemize}
			\item \code{np.eye(n)} creates and $n\times n$ identity matrix
			\item \code{A.T} trasponses the matrix \code{A}
			\item \code{np.linalg.norm(x,ord)} returns the norm (magnitude) of a vector \code{x} (optional argument \code{ord} for specifying the order of the norm)
			\item \code{np.matmul(A,B)} performs the matrix product of two vectors/matrices \\
					(there's also dot, inner, outer)
			\item \code{np.linalg.det(A)} returns the determinant of a matrix
			\item \code{np.linalg.inv(A)} returns the inverse of \code{A}
			\item \code{np.linalg.eig(A)} returns the eigenvalues and eigenvectors of a square matrix \\
					(can return complex numbers for non-symmetric matrices)
			\item \code{np.linalg.solve(A,b)} returns the solution of the system of linear equations $Ax=b$
		\end{itemize}
		
		A possible import alias: \\
		\code{from numpy import linalg as LA}
	\end{block}
	
\end{frame}

\begin{frame}	
	\frametitle{Linear Algebra Functions}
	
	\begin{exampleblock}{Solve a system of linear equations:}
		\footnotesize
		Do this WITHOUT and WITH \code{np.linalg.solve}
		\begin{itemize}
			\footnotesize
			\item Create a system of equations $Ax = b$ where $A$ is some matrix and $x,b$ are some vectors
			\item Solve the system by finding $x = A^{-1}b$
		\end{itemize}
	\end{exampleblock}
	
	\begin{exampleblock}{Find the principal components of some data:}
		\footnotesize
		\begin{itemize}
			\footnotesize
			\item Generate some random normal x data
			\item Generate y data that is linear w.r.t. x plus normal noise
			\item Evaluate correlation matrix (corrcoef)
			\item Evaluate eigenvectors and values of corrcoef
			\item Plot the vectors and the data
		\end{itemize}
	\end{exampleblock}
	
\end{frame}



% linear algebra, np.linalg
% eye, dot, det, inv





\end{document}




